\documentclass[12pt,a4paper]{article}

% ---------------------------------------------------------------
% Packages
% ---------------------------------------------------------------
\usepackage{amsmath,amssymb,amsthm}
\usepackage{geometry}
\geometry{margin=2.5cm}
\usepackage{natbib}
\bibliographystyle{plainnat}
\usepackage{hyperref}
\usepackage{booktabs}
\usepackage{array}
\usepackage{setspace}
\onehalfspacing

% ---------------------------------------------------------------
% Title
% ---------------------------------------------------------------
\title{\textbf{Multifractal Scaling of the CBOE Volatility Index:\\
First Application of the MMAR Framework to Implied Volatility\\
with Monte Carlo Validation}}

\author{[Author]\\
\small [Institution]\\
\small \texttt{[email]}}

\date{RSS International Conference, Bournemouth, 7--10 September 2026}

\begin{document}
\maketitle

% ---------------------------------------------------------------
% Abstract
% ---------------------------------------------------------------
\begin{abstract}
We present the first application of the Multifractal Model of Asset
Returns (MMAR; Mandelbrot, Fisher \& Calvet 1997) to an implied
volatility index. Using 8,136 daily observations of the CBOE VIX
spanning January 1990 to March 2026 --- a period encompassing four
economic cycles --- we implement the complete MMAR estimation pipeline:
partition functions, OLS-fitted scaling exponents, Legendre-transformed
multifractal spectrum, lognormal cascade calibration, fractional
Brownian motion generation, and Monte Carlo validation.

VIX log-returns strongly reject Gaussianity
($\text{KS } p < 10^{-10}$, excess kurtosis 6.40). The estimated
scaling function $\hat{\tau}(q) = -0.022q^2 + 0.207q - 1.019$ is
strictly concave, confirming genuine multiscaling. The Hurst exponent
$\hat{H} = 0.214$ places VIX in the strongly subdiffusive regime
(fractal dimension $d_H = 1.786$), consistent with its mean-reverting
character and in sharp contrast with the persistent scaling of equity
price indices. The most probable H\"{o}lder exponent is
$\hat{\alpha}_0 = 0.207$.

The estimated lognormal cascade variance $\hat{\sigma}^2 = -0.100$
is inadmissible under the standard MMAR parameterisation, providing
analytical evidence that the lognormal cascade is misspecified for
implied volatility. Formal validation via 10,000 Monte Carlo MMAR
simulations and 1,000 Gaussian benchmarks confirms this: MMAR produces
mean excess kurtosis of 0.07 against the empirical 6.40, with
Kolmogorov--Smirnov tests rejecting both MMAR and GBM at all
conventional significance levels.

We interpret the negative cascade variance and kurtosis failure as
joint evidence that VIX's multifractal structure belongs to a cascade
class outside the lognormal family --- a finding with direct
implications for volatility derivatives pricing and the modelling of
second-order financial time series.
\end{abstract}

\noindent\textbf{Keywords:} multifractal analysis, implied volatility,
Hurst exponent, partition functions, fractional Brownian motion,
Monte Carlo, VIX

\bigskip
\hrule
\bigskip

% ---------------------------------------------------------------
\section{Introduction}
% ---------------------------------------------------------------

Standard models of financial returns --- geometric Brownian motion (GBM)
and GARCH-family processes --- rest on assumptions that are
systematically violated by empirical data. Log-returns exhibit temporal
dependence, with periods of large changes clustering alongside periods
of small changes. Their distributions possess tails far heavier than
Gaussian and, crucially, these heavy tails do not vanish as the sampling
interval is extended: empirical studies confirm that multifractal
processes never converge to a Gaussian distribution regardless of
aggregation horizon. The CBOE Volatility Index (VIX), which distils
the market's expectation of 30-day forward volatility from S\&P 500
option prices, exemplifies these pathologies in a concentrated form:
excess kurtosis of 6.40, strong mean-reversion ($H \approx 0.25$ across
long lags), and regime-dependent jump behaviour during crises.

GARCH($p$,$q$) models \citep{bollerslev1986} address volatility
clustering by conditioning variance on past squared residuals, with
higher $\alpha$ coefficients producing sharper changes and $\beta$
capturing gradual dependence. Since $(\alpha + \beta) < 1$ in
stationary specifications, however, there is an inherent trade-off:
GARCH cannot simultaneously capture abrupt regime shifts and prolonged
volatility persistence. The multifractal framework resolves this by
constructing processes that are scale-invariant --- exhibiting identical
statistical properties at monthly, weekly, daily, or tick-level
frequencies --- with heavy tails arising endogenously from the cascade
structure rather than from distributional assumptions.

The Multifractal Model of Asset Returns (MMAR; \citealt{mandelbrot1997})
is the canonical continuous-time multifractal model for financial prices.
Over nearly three decades it has been applied to foreign exchange rates,
equity indices, individual stocks, commodities, and recently realised
variances. No published work, however, has applied the full MMAR
pipeline to any implied volatility index. This paper addresses that gap.
The contribution is threefold: (i) first MMAR application to the VIX
time series across 36 years of data; (ii) formal statistical validation
via Monte Carlo simulation and Kolmogorov--Smirnov tests, directly
addressing MMAR's well-known inference gap; and (iii) characterisation
of why the lognormal cascade is misspecified for VIX, with implications
for implied volatility modelling.

% ---------------------------------------------------------------
\section{Fractal Properties of Financial Time Series}
% ---------------------------------------------------------------

\subsection{Self-Similarity and the Hausdorff Dimension}

A fractal is a curve or geometric shape each part of which has the same
statistical properties as the whole --- exhibiting similar patterns at
progressively finer scales \citep{saupe1988}. For a one-dimensional
time series embedded in a two-dimensional plot, the roughness of the
trajectory is measured by the Hausdorff (fractal) dimension $d_H$,
which relates to the Hurst exponent via:
\begin{equation}
  d_H = 2 - H.
  \label{eq:hausdorff}
\end{equation}
When $d_H = 1$ the series is smooth, resembling a straight line; when
$d_H = 2$ it fills the plane entirely, like a jagged coastline. As $H$
increases the curve becomes smoother and $d_H$ decreases. For VIX with
$\hat{H} = 0.214$, equation~\eqref{eq:hausdorff} gives $d_H = 1.786$
--- a highly irregular, rough trajectory consistent with frequent
mean-reversions.

\subsection{Hurst Exponent and Diffusion Scaling}

The speed of diffusion of a time series is characterised by the mean
squared displacement (MSD):
\begin{equation}
  \operatorname{var}(\tau)
  = \bigl\langle |x_{t+\tau} - x_t|^2 \bigr\rangle.
  \label{eq:msd}
\end{equation}
For standard GBM, variance grows linearly with $\tau$ (normal
diffusion). Deviations from GBM follow a power-law:
\begin{equation}
  \operatorname{var}(\tau) \propto \tau^{2H}.
  \label{eq:var-hurst}
\end{equation}
For $H < \tfrac{1}{2}$: subdiffusion --- anti-persistent, mean-reverting,
MSD grows slower than linear. For $H > \tfrac{1}{2}$: superdiffusion ---
persistent, trending. For $H = \tfrac{1}{2}$: standard Brownian motion
with no memory. The fractal dimension $d_H = 2 - H$ maps the
persistence or anti-persistence of a time series to its geometric
roughness.

VIX exhibits $H \approx 0.249$ over long lags (1990--2022) and
$H \approx 0.398$ over short lags, both estimated via R/S analysis,
classifying it as subdiffusive. S\&P 500 returns, by contrast, show
$H \approx 0.637$ over long horizons --- reflecting the persistence of
equity price trends. Simulating a Hurst exponent alone is insufficient
to replicate VIX's character, however: fBm with any $H$ still produces
a Gaussian marginal distribution. What distinguishes VIX is the temporal
\emph{clustering} of volatility bursts --- days of extreme movement
followed by further extremes before abrupt subsidence. A leptokurtic
distribution generating occasional high-kurtosis realisations, without
this clustering, fails to reproduce the phenomenon.

\subsection{Long-Range Dependence}

Long-range dependent processes exhibit autocorrelation functions that
decay as power laws rather than exponentially:
\begin{equation}
  \rho(\tau \to \infty) \propto \tau^{-a}, \quad a = 2(1-H).
  \label{eq:acf-power}
\end{equation}
As $H \to 1$, $a \to 0$ and the decay becomes arbitrarily slow ---
indicating strong persistence. Fractionally integrated processes
\citep{robinson2003} provide one representation:
\begin{equation}
  (1 - L)^d S_t = \varepsilon_t, \quad d \notin \mathbb{Z},
  \label{eq:arfima}
\end{equation}
where the non-integer differencing parameter $d$ governs memory.
For VIX in the anti-persistent regime ($H < \tfrac{1}{2}$, $a > 1$),
autocorrelations of levels decay quickly, yet squared and absolute
returns display slow power-law decay --- a hallmark of multifractal
time series.

\subsection{Limitations of GARCH and GBM}

GBM is defined by $X(t) = X(0)\exp(\mu t + \sigma W(t))$ where $W(t)$
is standard Brownian motion. Its returns are lognormal, the process is
a martingale (memoryless), and the underlying is always positive.
These properties, while analytically convenient, make GBM a highly
simplified model of reality: it cannot produce fat tails, volatility
clustering, or long-range dependence. GARCH addresses the clustering
issue but cannot simultaneously model abrupt shifts and sustained long
cycles in volatility due to the stationarity constraint
$(\alpha + \beta) < 1$.

% ---------------------------------------------------------------
\section{The Multifractal Model of Asset Returns}
% ---------------------------------------------------------------

\subsection{Three Defining Features}

MMAR \citep{mandelbrot1997} unifies three properties absent from GARCH
and GBM:
\begin{enumerate}
  \item \textbf{Heavy tails with finite variance} --- arising
    endogenously from the multiplicative cascade structure.
  \item \textbf{Long-range dependence} --- via the fractional Brownian
    motion component \citep{mandelbrotness1968}.
  \item \textbf{Trading time} --- explicit modelling of the
    clock-time/natural-time relationship
    \citep{mandelbrottaylor1967}.
\end{enumerate}
Unlike L-stable distributions (which imply infinite variance) or fBm
alone (which retains Gaussian margins), MMAR combines long tails,
finite variance, and long dependence within a single framework.

\subsection{Three Assumptions: The Subordination Structure}

\textbf{Assumption 1.} Log-prices follow the compound process:
\begin{equation}
  X(t) \equiv B_H[\theta(t)],
  \label{eq:mmar-subordination}
\end{equation}
where $X(t) = \ln P(t) - \ln P(0)$, $B_H$ is fractional Brownian
motion with Hurst exponent $H$, and $\theta(t)$ is a stochastic trading
time. When plotted, fBm exhibits self-affinity: the plot looks the same
at any magnification.

\textbf{Assumption 2.} Trading time $\theta(t)$ is the cumulative
distribution function of a multifractal random measure. Not every day
``feels'' the same in financial markets: some days are fast (high
volume, intense price discovery) and others slow. This is the key
conceptual contribution of the MMAR. The multifractal measure captures
this heterogeneity hierarchically, across all time scales
simultaneously.

\textbf{Assumption 3.} $B_H$ and $\theta(t)$ are independent. Trading
time modifies \emph{when} fluctuations occur, not their direction or
correlation structure.

\subsection{Multiplicative Cascades}

Trading time is generated via a multiplicative cascade. At each step
$k$, the interval $[0,T]$ is divided into $b^k$ sub-intervals; each
receives a fraction of the total \emph{mass} (volatility intensity)
via an independent random multiplier $M_i$. Mass here is not
probability mass --- it is a scaling factor describing how volatility
intensity is distributed across time intervals, analogous to energy
distribution in turbulent flows.

Two cascade types arise. \emph{Conservative} cascades enforce
$\sum m_i = 1$ exactly at each step. \emph{Canonical} cascades require
only $\mathbb{E}[\sum M_i] = 1$ in expectation, allowing individual
multipliers to be drawn independently from a continuous distribution.
We employ a \textbf{canonical binomial cascade} ($b = 2$) with
lognormally distributed multipliers --- the parameterisation of
\citet{mandelbrot1997}.

Applying the cascade at progressively finer scales produces an uneven
mass distribution: some short intervals receive disproportionately high
intensity (volatility spikes), while neighbouring intervals receive
little. This hierarchical redistribution across scales generates the
clustering of high and low volatility periods that characterises
real financial time series --- and cannot be replicated by GARCH or GBM.

\subsection{Fractional Brownian Motion Component}

The fBm component $B_H(t)$ is defined by the Wiener integral:
\begin{equation}
  B_H(t) = \frac{1}{\Gamma(H+\tfrac{1}{2})}
  \left[
    \int_{-\infty}^0 \!\!\Bigl((t-s)^{H-\frac{1}{2}} - (-s)^{H-\frac{1}{2}}\Bigr)
    \,\mathrm{d}B(s)
    + \int_0^t (t-s)^{H-\frac{1}{2}} \,\mathrm{d}B(s)
  \right].
  \label{eq:fbm}
\end{equation}
The kernel weights $(t-s)^{H-1/2}$ introduce memory into the
increments. For $H < \tfrac{1}{2}$, the kernel decays quickly and
produces anti-persistent paths; for $H > \tfrac{1}{2}$, it decays
slowly, producing persistence. The fBm length is calibrated so that
the median standard deviation of simulated returns matches the empirical
value (0.0674 for VIX).

\subsection{MMAR Properties}

MMAR captures three empirical properties simultaneously:

\begin{table}[h!]
\centering
\begin{tabular}{ll}
\toprule
\textbf{Property} & \textbf{Mechanism} \\
\midrule
Non-Gaussian distribution (fat tails) & Multiplicative cascade heterogeneity \\
Heteroskedasticity ($H \neq \tfrac{1}{2}$) & Fractional Brownian motion \\
Volatility clustering & Trading time CDF maps fast/slow periods \\
\bottomrule
\end{tabular}
\caption{Three empirical properties captured by MMAR.}
\end{table}

% ---------------------------------------------------------------
\section{Estimation Methodology}
% ---------------------------------------------------------------

\subsection{Multifractality Condition}

A stochastic process $X(t)$ with stationary increments is
\emph{multifractal} if:
\begin{equation}
  \mathbb{E}\bigl[|X(t)|^q\bigr] = c(q)\,t^{\tau(q)+1},
  \quad \forall\, t \in T,\; q \in Q,
  \label{eq:multifractal-def}
\end{equation}
where $c(q)$ is a moment-dependent prefactor independent of $t$, and
$\tau(q)$ is the multifractal scaling exponent. The defining condition
for multifractality is \emph{non-linearity} of $\tau(q)$: a linear
$\tau(q)$ would imply a monofractal process (such as fBm) where a
single Hurst exponent governs all moments.

\subsection{Partition Function}

For moment orders $q \in [0.01, 30]$ (105 values) and time scales
$\Delta t$ that are integer factors of $T = 8136$ (24 scales), the
partition function is:
\begin{equation}
  S_q(T,\Delta t) = \sum_{i=0}^{N-1}
  \Bigl|X\!\bigl((i+1)\Delta t\bigr) - X\!\bigl(i\Delta t\bigr)\Bigr|^q,
  \label{eq:partition}
\end{equation}
where $N = T/\Delta t$. This is the sum of $q$-th power absolute
increments across all non-overlapping blocks of size $\Delta t$.

\subsection{Scaling Function}

OLS regression of $\ln S_q$ against $\ln \Delta t$ for each $q$ yields
\citep{calvet1997}:
\begin{equation}
  \ln\!\Bigl(\mathbb{E}\bigl[S_q(T,\Delta t)\bigr]\Bigr)
  = \hat{\tau}(q)\ln(\Delta t) + \hat{c}(q)\ln(T).
  \label{eq:ols-tau}
\end{equation}
The slope $\hat{\tau}(q)$ is the scaling exponent for moment order $q$.
Non-linearity of $\hat{\tau}(q)$ across $q$ confirms multiscaling. A
quadratic $\hat{\tau}(q) = aq^2 + bq + c$ with $a < 0$ is the
lognormal cascade signature; more general concave shapes indicate
other cascade families.

\subsection{Hurst Exponent}

$H$ is recovered by solving $\tau_X(1/H) = 0$ \citep{mandelbrot1997}:
\begin{equation}
  \hat{H} = \frac{1}{\hat{\tau}_X^{-1}(0)}.
  \label{eq:hurst-from-tau}
\end{equation}
This uses the root of the estimated $\hat{\tau}(q)$ curve, found by
fitting a quadratic and solving numerically. As a robustness check,
$H$ is also estimated via R/S analysis: $R/S(n) \propto n^H$, from
which $H$ is read off the log-log slope.

\subsection{Multifractal Spectrum via Legendre Transform}

The singularity spectrum $\hat{f}(\alpha)$ characterises the fractal
dimension of the set of time points with local H\"{o}lder exponent
$\alpha$. It is obtained via Legendre transform:
\begin{equation}
  \hat{f}(\alpha) = \inf_{q \in Q}\bigl\{q\alpha - \hat{\tau}(q)\bigr\} + D,
  \label{eq:legendre}
\end{equation}
where $D = 1$ is the embedding dimension. For a strictly concave
$\hat{\tau}(q)$, the infimum is achieved at the $q$ satisfying
$\hat{\tau}'(q) = \alpha$, yielding the classical conjugacy relation.

The spectrum answers: ``What is the fractal dimension of the set of
time points where the path has local H\"{o}lder exponent $\alpha$?''
Points of low $\alpha$ (rough, singular) correspond to volatility
spikes --- rare events, so $\hat{f}(\alpha_{\min}) < 1$. Points of
high $\alpha$ (smooth, regular) correspond to calm periods --- also
relatively rare. The most common behaviour, at $\hat{\alpha}_0$,
has $\hat{f}(\hat{\alpha}_0) = 1$, filling a set of full measure on
$[0,T]$. If the data were monofractal, $\hat{f}$ would degenerate to a
single point at $\alpha_0 = H$; the breadth of the estimated spectrum
quantifies the degree of multifractality.

The most probable H\"{o}lder exponent is:
\begin{equation}
  \hat{\alpha}_0 = \hat{f}^{-1}\!\bigl(\max_\alpha \hat{f}(\alpha)\bigr).
  \label{eq:alpha0}
\end{equation}
H\"{o}lder exponents should not be confused with the Hurst exponent:
$H$ is a global measure of diffusion scaling for the entire series;
$\alpha(t)$ are local measures of path regularity at specific points,
varying across time. They are related through the MMAR framework but
are conceptually distinct.

\subsection{Lognormal Cascade Parameters}

Given $\hat{H}$ and $\hat{\alpha}_0$, the canonical lognormal cascade
is parameterised by \citep{mandelbrot1997}:
\begin{equation}
  \hat{\lambda} = \frac{\hat{\alpha}_0}{\hat{H}},
  \qquad
  \hat{\sigma}^2 = \frac{2(\hat{\lambda}-1)}{\ln b}.
  \label{eq:cascade-params}
\end{equation}
These ensure the moment condition $\mathbb{E}[M] = b^{-1}$ for the
multiplier distribution $V = -\log_b M \sim \mathcal{N}(\hat{\lambda},
\hat{\sigma}^2)$:
\begin{equation}
  \mathbb{E}[M]
  = \exp\!\Bigl(-\hat{\lambda}\ln b
    + \tfrac{1}{2}\hat{\sigma}^2(\ln b)^2\Bigr)
  = \frac{1}{b}.
  \label{eq:moment-condition}
\end{equation}

\subsection{Trading Time and Full Simulation}

The trading time CDF is the cumulative sum of the cascade measure:
\begin{equation}
  \theta(t) = F_\mu(t) = \sum_{i=0}^t \mu(i).
  \label{eq:trading-time}
\end{equation}
The full MMAR simulation composes $X_t = B_H[\theta(t)]$, plugging
the trading time CDF into the fBm path to generate synthetic
log-price increments. The Python implementation adapts the MATLAB code
of \citet{wengert2010} for the cascade construction, using $K = 13$
dyadic levels ($2^{13} = 8192$ steps per path).

% ---------------------------------------------------------------
\section{Data}
% ---------------------------------------------------------------

The data consist of 8,136 daily closing observations of the CBOE
Volatility Index (VIX) from 2 January 1990 to 28 March 2026, sourced
from CBOE. This span covers four economic cycles: the 1990s dot-com
boom and bust, the 2000s housing boom and 2008 financial crisis, the
2010s cryptocurrency expansion, and the COVID-19 shock of 2020.
Covering multiple full cycles provides stable conditions for empirical
parameter estimation.

Descriptive statistics: mean price 19.53 (SD 8.01, price kurtosis
8.38). Log-returns $r_t = \ln(\text{VIX}_t) - \ln(\text{VIX}_{t-1})$
have mean $2.5 \times 10^{-5}$, standard deviation 0.0674, and excess
kurtosis 6.40. A Kolmogorov--Smirnov test for normality strongly
rejects Gaussianity ($p < 10^{-10}$), confirming fat tails as the
primary motivation for multifractal modelling.

% ---------------------------------------------------------------
\section{Results}
% ---------------------------------------------------------------

\subsection{Fractal Dimension and Diffusion Regime}

The estimated Hurst exponent $\hat{H} = 0.214$ implies fractal
dimension $d_H = 1.786$, placing VIX in the strongly subdiffusive
regime. This is consistent with R/S-based estimates from Section 2.2
(long-horizon $H \approx 0.249$). The near-identical values from two
independent estimation methods confirm the robustness of the
anti-persistence finding. The scaling
$\operatorname{var}(\tau) \propto \tau^{0.428}$ means VIX variance
grows at less than half the GBM rate.

\subsection{Multifractal Scaling Function and Spectrum}

The partition function $S_q(T,\Delta t)$ was computed for 105 moment
orders $q \in [0.01, 30]$ across 24 time scales. OLS regressions
confirm well-defined log-log linear relationships. The estimated
scaling function:
\begin{equation}
  \hat{\tau}(q) = -0.022q^2 + 0.207q - 1.019
  \label{eq:tau-estimated}
\end{equation}
is strictly concave (negative $q^2$ coefficient), confirming genuine
multifractality. At key moments: $\hat{\tau}(1) = -0.835$,
$\hat{\tau}(2) = -0.678$. The spectrum apex is at
$(\hat{\alpha}_0, \hat{f}(\hat{\alpha}_0)) = (0.207, 1.019)$.

\subsection{MMAR Parameter Estimates}

\begin{table}[h!]
\centering
\begin{tabular}{llp{7cm}}
\toprule
\textbf{Parameter} & \textbf{Value} & \textbf{Interpretation} \\
\midrule
$\hat{H}$ & 0.2143 & Strongly anti-persistent fBm \\
$\hat{\alpha}_0$ & 0.2069 & Most probable H\"{o}lder exponent \\
$\hat{\lambda}$ & 0.9654 & Cascade mean \\
$\hat{\sigma}^2$ & $-0.0997$ & \textbf{Inadmissible} under lognormal cascade \\
$b$ & 2 & Binomial cascade branching factor \\
\bottomrule
\end{tabular}
\caption{MMAR parameter estimates for VIX (8,136 daily obs., 1990--2026).}
\label{tab:params}
\end{table}

\subsection{Monte Carlo Validation}

10,000 MMAR paths ($K=13$, $2^{13} = 8192$ steps) and 1,000 GBM
benchmark paths were simulated with matched empirical mean and standard
deviation.

\begin{table}[h!]
\centering
\begin{tabular}{lrrr}
\toprule
\textbf{Metric} & \textbf{MMAR} & \textbf{GBM} & \textbf{Empirical VIX} \\
\midrule
Mean excess kurtosis & 0.068 & $-0.001$ & \textbf{6.40} \\
Mean log return & $1\times10^{-6}$ & $-3\times10^{-5}$ & $2.5\times10^{-5}$ \\
Mean SD of returns & 0.0958 & 0.0674 & 0.0674 \\
KS $p$-value (mean) & 0.000 & 0.000 & --- \\
\% paths KS $p > 0.05$ & 0.0\% & 0.0\% & --- \\
\bottomrule
\end{tabular}
\caption{Monte Carlo validation results: MMAR vs.\ GBM vs.\ empirical VIX.}
\label{tab:mc}
\end{table}

Both models are rejected at all conventional significance levels.
The MMAR kurtosis shortfall is stark: lognormal cascade simulations
generate mean kurtosis of 0.07 against an empirical value of 6.40 ---
a 95-fold underestimate, replicated consistently across 10,000 paths.
For VIX, the cascade fails not by generating too little tail mass
(its typical failure mode for equity returns), but by being structurally
unable to concentrate tail mass at all given the inadmissible parameter
regime.

% ---------------------------------------------------------------
\section{Discussion}
% ---------------------------------------------------------------

\subsection{The $\hat{\sigma}^2 < 0$ Finding: Cascade Misspecification}

The inadmissible cascade variance is the paper's most theoretically
consequential result. From the parameterisation
\eqref{eq:cascade-params}, $\hat{\sigma}^2 > 0$ requires
$\hat{\lambda} > 1$, equivalently $\hat{\alpha}_0 > \hat{H}$. For
equity returns this holds comfortably: S\&P 500 estimates yield
$H \approx 0.60$, $\alpha_0 \approx 0.65$, giving $\lambda \approx
1.08$ and $\sigma^2 \approx 0.12$. For VIX, $\hat{\lambda} =
0.207/0.214 = 0.965 < 1$. This is a structural consequence of the
strongly anti-persistent Hurst exponent, not a numerical artefact.

The condition $\hat{\alpha}_0 < \hat{H}$ has a concrete
interpretation. In the MMAR framework, $\alpha_0$ is the most probable
local H\"{o}lder exponent of the compounded process
$X(t) = B_H[\theta(t)]$. For a lognormal cascade to be identifiable,
the trading time deformation must \emph{amplify} local irregularity ---
stretching some intervals and compressing others to produce more
variation than the raw fBm. When $\hat{\alpha}_0 < \hat{H}$, the data
imply the opposite: trading time deformation would need to
\emph{smooth} local fluctuations relative to fBm. No lognormal
multiplier distribution satisfying the conservation constraint
$\mathbb{E}[M] = 1/b$ can achieve this. The lognormal family's
parameter boundary has been reached.

Three interpretations are plausible. First, VIX may require a
sub-lognormal cascade --- one with more concentrated multiplier
distribution and less mass redistribution per cascade step.
Beta-distributed multipliers on $[0,1]$ naturally enforce this
concentration. Second, VIX's mean-reversion may be incompatible with
Assumption 3 (independence of $B_H$ and $\theta(t)$): when volatility
spikes, both trading intensity and subsequent price direction are
affected, violating independence. Third, the anti-persistent fBm with
$H = 0.214$ may itself be incorrect; rough volatility models
\citep[Gatheral et al.\ 2018]{} use $H \approx 0.1$ for variance
processes, and a lower $H$ could shift $\hat{\alpha}_0/\hat{H}$ above
unity and restore admissibility.

\subsection{GARCH, MSM, and the MMAR Framework}

The GARCH trade-off --- between capturing abrupt shifts (high $\alpha$)
and sustained cycles (high $\beta$) --- is an inherent constraint of
the stationarity condition $(\alpha + \beta) < 1$. MMAR's failure is
different in character: the framework is conceptually correct for VIX
in that it provides a mechanism for time-varying variance through
trading time, but the \emph{lognormal} cascade specification is
incorrectly calibrated for the subdiffusive regime in which VIX operates.

\citet{luxliu2017} and \citet{liu2019} showed that Markov-Switching
Multifractal (MSM) models outperform GARCH at long horizons, and
GARCH outperforms at short horizons. Neither study applied MSM to VIX
as a target variable. The present finding that MMAR is misspecified at
the continuous-time level raises the question of whether MSM's
discrete cascade structure --- with bounded, finitely many volatility
components --- would fare better for implied volatility, given VIX's
bounded mean-reversion dynamics.

\subsection{Implications for Alternative Cascade Specifications}

The negative $\hat{\sigma}^2$ and kurtosis failure together indicate
that VIX's multifractal trading time does not follow a lognormal
cascade. Candidates for replacement include:
\begin{itemize}
  \item \textbf{Poisson cascades}: sparse, jump-like intensity
    redistributions, appropriate for the discrete crisis-onset
    character of VIX spikes;
  \item \textbf{$\alpha$-stable cascades}: permitting heavy-tailed
    multiplier distributions with infinite variance;
  \item \textbf{Beta-distributed cascades}: bounded multipliers
    satisfying the conservation constraint while allowing
    sub-lognormal concentration.
\end{itemize}
Each generates a different $\tau(q)$ profile and requires different
parameter identification strategies. Comparing these specifications
against the empirical VIX spectrum is a direct extension of this work.

% ---------------------------------------------------------------
\section{Conclusion}
% ---------------------------------------------------------------

This paper presents the first application of the complete MMAR pipeline
to an implied volatility index. Using 8,136 daily VIX observations from
January 1990, we establish five results:
\begin{enumerate}
  \item VIX log-returns are non-Gaussian with excess kurtosis 6.40,
    confirmed by KS test ($p < 10^{-10}$).
  \item VIX is genuinely multifractal: the scaling function
    $\hat{\tau}(q) = -0.022q^2 + 0.207q - 1.019$ is strictly concave.
  \item $\hat{H} = 0.214$ places VIX in the strongly subdiffusive,
    anti-persistent regime ($d_H = 1.786$), in sharp contrast with
    equity price indices.
  \item The lognormal cascade variance $\hat{\sigma}^2 = -0.100$ is
    inadmissible, indicating structural cascade misspecification.
  \item 10,000 MMAR and 1,000 GBM simulations confirm that neither
    model replicates VIX's tail behaviour (MMAR kurtosis 0.07 vs.\
    empirical 6.40; KS rejection at 100\% of simulated paths).
\end{enumerate}

These results have direct implications for the \$50+ billion VIX
derivatives ecosystem. Systematic underestimation of kurtosis by
standard models suggests the field requires alternative cascade
specifications for implied volatility. Future work will explore Poisson
and $\alpha$-stable cascade alternatives and extend the analysis to
the VVIX (volatility-of-volatility index), which represents the
second-order derivative of the quantity studied here.

% ---------------------------------------------------------------
% Bibliography
% ---------------------------------------------------------------
\bibliography{RSS2026_MMAR_VIX}

\end{document}
